\section{A notation map, with places to return to}\label{appE:sec:notation}
The table below lists the main symbols of the book in the order in which they appear, together with the place where each one is explained. Use it to find your way back to an explanation. When a symbol in a later chapter looks unfamiliar, the one-line summary here can remind you what it means, but rereading the section where it was introduced, with its examples, is what makes it usable again.

A few symbols do more than one job, and the surrounding sentence decides which job is meant.
\begin{itemize}
\item Vertical bars: for a real number $x$, $|x|$ is its absolute value; for a finite set $S$, $|S|$ is the number of its members.
\item Brackets: $(1,3)$ can be the interval of real numbers strictly between $1$ and $3$, or the ordered pair with first entry $1$ and second entry $3$. Square brackets also have two uses: $[1,3]$ is an interval, while $[x]$ is the equivalence class of $x$.
\item The arrow $\to$: in $f:A\to B$ it names the input and output sets of a function; in $u\to v$ it is a directed edge of a graph; in $x_n\to x_*$ it says that a sequence converges. The double arrow $\Rightarrow$ is different again: it joins two statements.
\item The notation $f^{-1}(T)$ is a preimage. It makes sense for every function $f$, whether or not $f$ has an inverse.
\item A difference of sets such as $A\setminus B$ is a set, not a number. It removes members; it does not subtract values.
\item Letters are reused from chapter to chapter. For example, $q$ is a quotient in Chapter~\ref{ch:induction} and a contraction factor in Chapters~\ref{ch:contraction} and~\ref{ch:capstone}; $I$ is an identity matrix in Chapter~\ref{ch:linear}, while $I_A$ is the indicator of an event in Chapter~\ref{ch:probability}.
\end{itemize}

\begin{longtable}{@{}p{1.25in}p{3.95in}@{}}
\toprule Symbol & Meaning, and where it is explained\\\midrule\endhead
$\NN,\ZZ,\QQ,\RR$ & The natural numbers $0,1,2,\ldots$ (zero included), the integers, the rational numbers, and the real numbers (Section~\ref{ch01:sec:notation}).\\
$\sum_{j=a}^{b}x_j$ & The sum $x_a+x_{a+1}+\cdots+x_b$, with one term for each integer $j$ from $a$ to $b$ (Section~\ref{ch01:sec:notation}). When $b<a$ there are no terms, and this \emph{empty sum} is $0$ (Section~\ref{ch05:sec:first-induction}). A sum written $\sum_{\omega\in\Omega}$ has one term for each member of the finite set $\Omega$ (Section~\ref{ch13:sec:model}).\\
$x<y$, $x\le y$ & $x$ is strictly smaller than $y$; $x$ is smaller than or equal to $y$. The two differ only when $x=y$ (Section~\ref{ch01:sec:notation}).\\
$\sqrt t$ & For a real number $t\ge0$, the nonnegative number whose square is $t$. So $\sqrt9=3$, not $-3$ (Section~\ref{ch01:sec:notation}).\\
$P\Rightarrow Q$, $P\Leftrightarrow Q$ & ``If $P$, then $Q$''; ``$P$ if and only if $Q$,'' which means both $P\Rightarrow Q$ and $Q\Rightarrow P$. Writing either one states a claim; it does not say that the claim has been proved (Section~\ref{ch02:sec:implication}).\\
$\neg,\land,\lor$ & Not; and; or. The ``or'' is inclusive: $P\lor Q$ is true when at least one of $P$ and $Q$ is true, including when both are (Section~\ref{ch02:sec:connectives}).\\
$\forall,\exists$ & For every; there exists (at least one). Quantifiers are read from left to right, and a value chosen at an $\exists$ may depend on every value chosen before it, but not on anything chosen after it. So changing the order can change the meaning (Section~\ref{ch02:sec:quantifiers}).\\
$\eps,\delta$ & Traditional names for tolerances, that is, positive amounts of error we are prepared to allow. They are names, not fixed numbers: ``for every $\eps>0$'' covers $\eps=1/10$, $\eps=1/1000$, and every other positive value (Sections~\ref{ch02:sec:quantifiers} and~\ref{ch11:sec:distance}). In Chapter~\ref{ch:capstone}, $\eps_n$ instead bounds the error made in the update at step~$n$, the one that produces $y_{n+1}$.\\
$|x|$ & The absolute value of a real number $x$: its distance from $0$ on the number line (Section~\ref{sec:magnitude}).\\
$[a,b]$, $(a,b)$, $[0,\infty)$ & Intervals of real numbers. A square bracket includes its endpoint and a round bracket excludes it; $[0,\infty)$ is the set of all real $x$ with $x\ge0$ (Section~\ref{sec:magnitude}).\\
$x\in A$, $x\notin A$ & $x$ is a member of the set $A$; $x$ is not a member of $A$ (Section~\ref{ch03:sec:sets}).\\
$A\subseteq B$, $\varnothing$ & Every member of $A$ is also a member of $B$; the empty set, which has no members (Section~\ref{ch03:sec:sets}).\\
$\{x\in A:P(x)\}$ & The set of members of $A$ that have the property $P$. Read the colon as ``such that'' (Section~\ref{ch03:sec:sets}).\\
$A\cup B$, $A\cap B$ & Union: the objects in at least one of $A$ and $B$. Intersection: the objects in both (Section~\ref{ch03:sec:sets}).\\
$A\setminus B$ & The members of $A$ that are not in $B$ (Section~\ref{ch03:sec:sets}).\\
$A\times B$ & The set of ordered pairs $(a,b)$ with $a\in A$ and $b\in B$ (Section~\ref{ch03:sec:sets}).\\
$|S|$ & The number of members of a finite set $S$, each counted once; $|\varnothing|=0$ (Section~\ref{ch03:sec:sets}).\\
$f:A\to B$, $x\mapsto x^2$ & A function with domain $A$ and codomain $B$; the function that sends each input $x$ to $x^2$ (Section~\ref{ch03:sec:functions}).\\
$g\circ f$, $\id_A$ & ``$g$ after $f$'': apply $f$ first, then $g$. The identity function on $A$, which sends every element to itself (Section~\ref{ch03:sec:functions}).\\
$f(S)$, $f^{-1}(T)$ & The image of a set $S$ of inputs: all outputs $f(x)$ with $x\in S$. The preimage of a set $T$ of outputs: all inputs $x$ with $f(x)\in T$ (Chapter~\ref{ch:maps}).\\
$a\bmod m$ & The remainder when the integer $a$ is divided by the positive integer $m$; it always lies between $0$ and $m-1$ (Section~\ref{ch05:sec:euclid}).\\
$x\sim y$, $[x]$ & $x$ is related to $y$; the equivalence class of $x$, the set of all elements related to $x$ (Section~\ref{ch06:sec:differences}).\\
$a\equiv b\pmod m$ & $a$ is congruent to $b$ modulo $m$: the positive integer $m$ divides $a-b$ (Section~\ref{ch06:sec:residues}).\\
$[a]_m$, $\ZZ/m\ZZ$ & The residue class of $a$: the whole set of integers congruent to $a$ modulo $m$, not just the one integer $a$. The set of all $m$ residue classes (Section~\ref{ch06:sec:residues}).\\
$x\preceq y$ & $x$ is below $y$ in a partial order (Section~\ref{ch06:sec:order}).\\
$\deg(v)$ & The degree of a vertex $v$: the number of its neighbors (Section~\ref{ch07:sec:graph}).\\
$u\to v$ & A directed edge from $u$ to $v$; $v$ is then an out-neighbor of $u$ (Section~\ref{ch07:sec:sink}).\\
$N_1^+(v)$, $N_2^+(v)$ & The first out-neighborhood of $v$: the vertices at directed distance exactly one from $v$. The second out-neighborhood: the vertices at directed distance exactly two (Section~\ref{ch07:sec:seymour}).\\
$N(s)$, $N(S)$ & The tasks accepted by the student $s$; the tasks accepted by at least one student of the set $S$ (Sections~\ref{ch08:sec:model} and~\ref{ch08:sec:shortage}).\\
$\RR^n$ & The set of vectors $(x_1,\ldots,x_n)$ with $n$ real coordinates (Section~\ref{ch10:sec:vectors}).\\
$A_{ij}$, $Ax$ & The entry of the matrix $A$ in row $i$ and column $j$; the vector obtained by applying $A$ to $x$ (Section~\ref{ch10:sec:linearmaps}).\\
$\ker T$ & The kernel of a linear map $T$: the inputs that $T$ sends to the zero vector (Section~\ref{ch10:sec:kernel}).\\
$u\cdot v$ & The dot product: multiply corresponding coordinates and add. The result is a number, not a vector (Section~\ref{ch10:sec:inner}).\\
$\norm{x}_2$ & The Euclidean length $\sqrt{x\cdot x}$ of a vector. The subscript $2$ is part of the name, not a coordinate (Section~\ref{ch10:sec:inner}).\\
$A^{\mathsf T}$, $I$ & The transpose of $A$, obtained by exchanging rows and columns. The identity matrix, with ones on the diagonal and zeros elsewhere (Section~\ref{ch10:sec:product}).\\
$(X,d)$, $d(x,y)$ & A metric space: a set $X$ with a chosen metric $d$; the distance from $x$ to $y$. The choice of metric is part of the problem (Section~\ref{ch11:sec:metric}).\\
$d_2$, $d_\infty$ & Two metrics on $\RR^n$: the Euclidean distance $\norm{x-y}_2$, and the largest of the coordinate differences $|x_j-y_j|$ (Section~\ref{ch11:sec:metric}).\\
$x_n\to x_*$ & The sequence converges to $x_*$: for every tolerance $\eps>0$, all terms from some index on are within distance $\eps$ of $x_*$ (Section~\ref{ch11:sec:convergence}).\\
$x_*$ & A name for a special point: a limit in Chapter~\ref{ch:limits}, a fixed point in Chapters~\ref{ch:contraction} and~\ref{ch:capstone}. The star is part of the name.\\
$Tx$ & Short for $T(x)$ (Section~\ref{ch12:sec:iteration}).\\
$\Omega$, $p(\omega)$ & The sample space, a finite set of outcomes; the probability of the outcome $\omega$ (Section~\ref{ch13:sec:model}).\\
$\Prob(A)$ & The probability of an event $A\subseteq\Omega$: the sum of $p(\omega)$ over the outcomes $\omega\in A$ (Section~\ref{ch13:sec:model}).\\
$\E X$ & The expectation of a random variable $X$: the weighted average $\sum_{\omega\in\Omega}p(\omega)X(\omega)$ (Section~\ref{ch13:sec:expectation}).\\
$I_A$ & The indicator of the event $A$: it equals $1$ at outcomes in $A$ and $0$ at all other outcomes (Section~\ref{ch13:sec:indicators}).\\
$A_s$, $C_s$ & The aperiodic autocorrelation at shift $s$, defined for $1\le s\le n-1$, a sum of $n-s$ products in which nothing wraps around. The periodic autocorrelation at shift $s$, defined for $0\le s\le n-1$, a sum of $n$ products in which positions are taken modulo $n$ (Sections~\ref{ch14:sec:aperiodic} and~\ref{ch14:sec:periodic}).\\
\code{a : T}, \code{f x}, \code{:=} & In Lean: \code{a} has type \code{T} (for a proposition \code{P}, \code{h : P} says that \code{h} is a proof of \code{P}); the function \code{f} applied to \code{x}; ``is defined as'' or ``is proved by'' (Section~\ref{ch15:sec:symbols}).\\
\code{sorry} & In Lean, a placeholder for a missing proof. The command \code{\#print axioms} shows whether a theorem depends on one (Section~\ref{ch16:sec:axioms}).\\
$B(c,r)$ & The ball with center $c$ and radius $r>0$: the points whose distance from $c$ is strictly less than $r$ (Section~\ref{ch17:sec:balls}).\\
\bottomrule
\end{longtable}
